% 真矢量重建：沿用嵌套集合、样本约束与算法选点的已核验模型设计。
% 插值集合/点位置没有数量或风险数值含义；训练样本直接进入A。
\begin{tikzpicture}[foundation picture]
 \path[use as bounding box] (0,0) rectangle (112,72);
 \draw[FigureStroke,line width=1.1pt,fill=white] (32,31) ellipse (31 and 28);
 \node (architecture) at (9,66) {架构};
 \draw[FigureStroke,line width=.9pt] (architecture.south)--(11,51.6);
 \node at (17,48) {$\mathcal H$};
 \node at (17,42) {假设空间};
 \draw[FigureStroke,line width=1.1pt,fill=FigureBlueFill]
  (22,22)..controls(20,30)and(29,34)..(35,35)
  ..controls(44,40)and(46,34)..(51,32)
  ..controls(58,30)and(60,23)..(55,19)
  ..controls(49,13)and(39,14)..(32,13)
  ..controls(25,12)and(20,17)..(22,22)--cycle;
 \foreach \x/\y in {
  8/33,12/20,13/40,17/13,20/25,21/54,27/20,28/51,31/43,35/48,
  41/52,43/41,49/44,56/39,58/30,53/14,48/10,43/6,35/7,26/9,19/37,
  27/26,32/18,44/17,54/18}{
  \fill[FigureStroke] (\x,\y) circle (.7);
 }
 \node at (40,31) {插值解};
 \node at (40,26) {$\hat R_S(h)=0$};
 \fill[FigureRed] (52,20) circle (1.4);
 \node at (40,20) {算法输出};
 \node[foundation node,fill=white,minimum width=36mm,minimum height=14mm,align=center] (sample)
  at (90,62) {训练样本$S$\\$(x_1,y_1),\ldots,(x_n,y_n)$};
 \draw[FigureMuted,dashed,line width=.8pt] (sample.south west)--(49,36);
 \node[fill=white,inner sep=1pt] at (64,45) {约束};
 \foreach \x/\label in {77/初始化,91/优化器,105/正则化}{
  \node[foundation node,fill=white,minimum width=10mm,minimum height=7mm] (config\x)
   at (\x,39) {\label};
 }
 \node[foundation model,fill=white,minimum width=34mm,minimum height=11mm] (algorithm)
  at (92,20) {算法$A$};
 \foreach \x in {77,91,105}{
  \draw[foundation flow] (config\x.south)--(\x,25.5);
 }
 \draw[foundation flow] (sample.east)--(111,62)--(111,20)--(algorithm.east);
 \draw[foundation flow,draw=FigureRed] (algorithm.west)--(53.6,20);
\end{tikzpicture}%
